\scanpage{37}
\mspage{36}
\mspage{83}
\margin{n° 382}

des germes en x de sections qui ont un contact d'ordre $\geq k'$ avec la section nulle ; cet espace est la fibre en x d'un sous-fibré vectoriel de $P^k(X,E)$, qui est aussi un module sur $P^k(X)$. En particulier, le noyau de $P^k(X,E)\longrightarrow P^{k-1}(X,E)$ s'identifie au produit tensoriel $S^k(T(X)^*)\otimes E$ de E avec la k-ième puissance symétrique du dual du fibré tangent à X.

7.2.4. Soient X, Y des variétés de classe $C^r$, $\varpi:Y\longrightarrow X$ une submersion, $k'\leq k$ des entiers finis $\leq r$, $\ell=k-k'$, et soit $\varpi_1$ la projection de $P^{k'}(X,\varpi)$ sur X. L'application canonique de $J^k(X,Y)$ dans $J^\ell(X,J^{k'}(X,Y))$ définie au n° 7.1.10 induit une application de classe $C^{r-k}$ de $P^k(X,\varpi)$ dans $P^\ell(X,\varpi_1)$ qui, lorsque K est de caractéristique zéro, est un isomorphisme de $P^k(X,\varpi)$ sur une sous-variété de $P^\ell(X,\varpi_1)$ ; cette application est un morphisme de fibrés sur X si X et Y sont des variétés pures.

7.2.5. Soient X, Y, Y' des variétés de classe $C^r$, $\varpi:Y\longrightarrow X$ et $\varpi':Y'\longrightarrow X$ des submersions, $g:Y\longrightarrow Y'$ un morphisme de Y dans Y' tel que $\varpi'\circ g=\varpi$, \inserted{(1) et k un entier $\leq r$. L'application $j^k(\mathrm{id}.,g)$} (cf. n° 7.1.11) induit une application $P^k(g):P^k(X,\varpi)\longrightarrow P^k(X,\varpi')$ de classe $C^{r-k}$ qui, lorsque X, Y, Y' sont des variétés pures, est un morphisme pour les fibrations sur X. Si $Y''$ est une variété de classe $C^r$, $\varpi'':Y''\longrightarrow X$ une submersion et $g':Y'\longrightarrow Y''$ un morphisme tel que $\varpi'=\varpi''\circ g'$, alors $P^k(g'\circ g)=P^k(g')\circ P^k(g)$. Si $E\xrightarrow{\lambda}F\xrightarrow{\mu}G$ est une suite exacte de fibrés vectoriels sur X, alors
\diagram{\begin{tikzcd}
P^k(X,E) \arrow[r,"P^k(\lambda)"] & P^k(X,F) \arrow[r,"P^k(\mu)"] & P^k(X,G)
\end{tikzcd}}
est une suite exacte. Si A et B sont des fibrés vectoriels sur X, alors
\[
P_x^k(X,A\otimes B)\simeq P_x^k(X,A)\otimes_{P_x^k(X)}P_x^k(X,B),\quad\text{pour tout }x\in X.
\]

7.2.6. Soient X, X', Y' des variétés de classe $C^r$, $f:X\longrightarrow X'$ un morphisme, $\varpi:Y'\longrightarrow X'$ une submersion, $Y=X\times_{X'}Y'$ (n° 4.11.3) et $\varpi$ la submersion canonique de Y sur X, $(\qquad)$. L'application $j^k(f,\mathrm{Id}.)$ (cf. n° 7.1.11) induit une
